\documentclass[12pt]{article}

% --- Paquetes Esenciales ---
\usepackage[utf8]{inputenc}    % Caracteres en español
\usepackage[spanish]{babel}   % Idioma español
\usepackage{amsmath}          % Fórmulas matemáticas
\usepackage{geometry}         % Ajuste de márgenes
\geometry{letterpaper, margin=1in}

% --- Paquete para Circuitos ---
\usepackage[siunitx, american]{circuitikz}

\begin{document}

\section{Ejercicio de realizar un circuito}

\begin{center}
\begin{circuitikz}[scale=1.4]
    % --- LINEAS PRINCIPALES DEL BUS (SUPERIOR E INFERIOR) ---
    \draw (0,4) -- (6,4);
    \draw (0,0) -- (6,0);
    
    % --- RAMA 1: FUENTE DE VOLTAJE DC (12V) ---
    \draw (0,0) to[battery, l=$12\text{V}$] (0,4);
    % Signos de polaridad manuales para asegurar que salgan bien
    \node at (0.3, 2.3) {$+$};
    \node at (0.3, 1.7) {$-$};
          
    % --- RAMA 2: RESISTENCIA R1 ---
    \draw (2,4) to[R] (2,0);
    % Etiquetas al lado de R1
    \node[right] at (2, 2.2) {$R_1 = 2\,\Omega$};
    \node[right] at (2, 2.5) {$V_1$};
    \node[left] at (2, 1.2) {$I_1$};
          
    % --- RAMA 3: RESISTENCIA R2 ---
    \draw (4,4) to[R] (4,0);
    % Etiquetas al lado de R2
    \node[right] at (4, 2.2) {$R_2 = 7\,\Omega$};
    \node[right] at (4, 2.5) {$V_2$};
    \node[left] at (4, 1.2) {$I_2$};
          
    % --- RAMA 4: RESISTENCIA R3 ---
    \draw (6,4) to[R] (6,0);
    % Etiquetas al lado de R3
    \node[right] at (6, 2.2) {$R_3 = 4\,\Omega$};
    \node[right] at (6, 2.5) {$V_3$};
    \node[left] at (6, 1.2) {$I_3$};

\end{circuitikz}
\end{center}

\end{document}